\documentclass[11pt,a4paper]{article}
\usepackage[margin=27mm]{geometry}
\usepackage{amsmath,amssymb,amsthm,hyperref}
\newtheorem{proposition}{Proposition}
\newcommand{\Z}{\mathbb Z}
\newcommand{\one}{\mathbf 1}
\title{The rank-four toroidal families in cubic coordinates}
\author{Human Compression project}
\date{}
\begin{document}\maketitle

The fixed manuscript~\cite{FLW} displays three period families for each
of the types $\widetilde B_3$ and $\widetilde A_3$ on page~11.
They follow directly from the acquired cubic halving construction.
We give the exact change of parameters, including the smallest allowed
periods. This proves existence and the marked group descriptions of the
displayed families. Their exhaustion is the classification attributed in
the source to Ens; it is not a consequence of this parameter calculation.

\section{The acquired construction}
Put
\[
 D=\{v\in\Z^3:v_1+v_2+v_3\equiv0\pmod2\},\qquad
 E=2\Z^3+\Z\one.
\]
Write $B_3$ for the signed permutation group, and $B_3^+$ for the
subgroup with an even number of sign changes. Let $z_i$ change the sign
of coordinate $i$, let $s_i$ interchange coordinates $i,i+1$, and let
$\tau_v$ be translation by $v$. The acquired affine calculation gives
\[
 x=\tau_{(1,1,0)}z_1z_2s_1,\qquad y=z_2z_3s_2,
\]
\[
 \langle x,s_1,s_2,z_3\rangle=D\rtimes B_3,\qquad
 \langle x,s_1,s_2,y\rangle=D\rtimes B_3^+.
\]
Their diagrams have, respectively, the bonds $02,12$ labelled three
and $23$ labelled four, and the four-cycle $02,12,13,03$ with all
labels three. Thus these are the marked affine groups of types
$\widetilde B_3$ and $\widetilde A_3$.

For $t\ge2$ put
\[
 \Lambda_1(t)=t\Z^3,\qquad \Lambda_2(t)=tD,\qquad
 \Lambda_3(t)=tE,\qquad L_k(t)=D\cap\Lambda_k(t).
\]
The acquired geometry proof applies one halving at the left end of the
cubic toroid, and then one at the opposite end. It gives actual regular
hypertopes $H_{10}$ and $H_{11}$, with the above markings and groups
\[
 (D/L_k(t))\rtimes B_3,\qquad
 (D/L_k(t))\rtimes B_3^+,
\]
for every $k\in\{1,2,3\}$ and $t\ge2$ except $(k,t)=(1,2)$.
Here existence includes thinness, residual connectivity and chamber
transitivity, proved in the acquired incidence models. These conclusions
are not being inferred from the diagrams or group orders. In particular,
the case $(k,t)=(2,2)$ is allowed for $H_{10}$ and $H_{11}$: its
separate exclusion concerns a second halving at the same end.

\section{The source parameters}
The period generated by a source vector means the integer span of its
orbit under the corresponding finite linear group. Both finite groups
above give the same three spans in dimension three.

\begin{proposition}
The following table identifies every displayed source family with the
acquired construction. The orders in the last two columns are the orders
of the full type-preserving groups.
\[
\begin{array}{c|c|c|c|c}
 \text{source vector and domain}&(k,t)&L_k(t)&
 |G_{\widetilde B_3}|&|G_{\widetilde A_3}|\\\hline
 (2s,0,0),\ s\ge2&(1,2s)&2s\Z^3&192s^3&96s^3\\
 (s,s,0),\ s\ge2&(2,s)&sD&48s^3&24s^3\\
 (2s,2s,2s),\ s\ge1&(3,2s)&2sE&768s^3&384s^3
\end{array}
\]
The source objects are the corresponding incidence quotients of the
universal marked affine complexes.
\end{proposition}
\begin{proof}
The orbit of $e_1$ under either group is the set of signed coordinate
vectors: an unwanted sign change can be placed in a zero coordinate.
Its span is $\Z^3$. Similarly the orbit of $e_1+e_2$ consists of all
roots $\pm e_i\pm e_j$. These span $D$: for example
$e_1+e_2,e_2-e_1,e_3-e_2$ are a lattice basis of determinant two.

Under $B_3$, the orbit of $\one$ spans $E$, since sign differences
give $2e_i$. Under $B_3^+$ it spans
$M=2D+\Z\one$: differences from $\one$ are twice the
characteristic vectors of two-element subsets, whose span is $D$.
In dimension three $M=E$. Indeed $\one-e_i\in D$, and
$2e_i=2\one-2(\one-e_i)\in M$. The reverse inclusion is immediate.
Scaling gives precisely the three source period lattices in the table.

Each of these lattices lies in $D$. The first and third do so because
their coordinates are even, and the middle does by definition.
Consequently their intersections with $D$ are unchanged. Substitution
in $L_k(t)$ proves the parameter correspondence. The first row has
$t\ge4$, the second has $t\ge2$ with $k=2$, and the third has
$t\ge2$ with $k=3$. None is the excluded pair $(1,2)$, so both
acquired geometry constructions apply throughout the exact source
domains, including $s=1$ in the third row.

Finally $[\Z^3:D]=2$, $[\Z^3:E]=4$, and scaling a lattice by $a$
multiplies its index by $a^3$. Thus the three indices in $D$ are
$4s^3,s^3,16s^3$. Multiply these by $|B_3|=48$ or
$|B_3^+|=24$. The acquired coset identification of each actual
geometry gives the stated incidence quotient, with its actual marking.
\end{proof}

\paragraph{Attribution and scope.}
The source attributes the complete rank-four family classification to
Ens. The cubic toroid input is the retained McMullen--Schulte
classification, and the halving constructions originate with the earlier
authors cited in the acquired proof. The present addition is the literal
parameter and marking interface between those proofs and this source.
It does not prove that every invariant lattice is in the table, relax a
small-period condition, or claim a new conjecture answer. No new finite
census, Lean formalization or compression measurement is involved.

\begin{thebibliography}{9}
\bibitem{FLW} M.~E.~Fernandes, D.~Leemans and A.~Ivi\'c Weiss,
\emph{An exploration of locally spherical regular hypertopes},
fixed catalogue L033 manuscript, pages 10--11.
\bibitem{acquired} Human Compression project,
\emph{Generalized halving: corrected source recovery},
accepted affine group calculation and iterated geometry proof,
Sections 20 and 30 in the frozen application sources.
\end{thebibliography}
\end{document}
